\begin{table}[ht]
\centering
\footnotesize
\caption{Robustness: Effect on Jail Population Rate}
\label{tab:robustness}
\begin{tabular}{@{}l@{\hskip 6pt}c@{\hskip 6pt}c@{\hskip 6pt}c@{\hskip 6pt}c@{\hskip 6pt}c@{\hskip 6pt}c@{\hskip 6pt}c@{}}
\toprule
 & (1) & (2) & (3) & (4) & (5) & (6) & (7) \\
 & Full & Pre-COVID & Pre-2020 & No 2020 & Pop-Wt & AAPI & Donut \\
\midrule
Progressive DA & $-$179.1*** & $-$217.3*** & $-$202.9*** & $-$183.8*** & $-$54.5*** & 5.4 & $-$184.9*** \\
 & (33.0) & (47.3) & (40.5) & (32.5) & (17.6) & (5.9) & (33.3) \\
\midrule
$N$ & 52,704 & 44,116 & 52,495 & 50,428 & 52,583 & 17,087 & 50,704 \\
$R^2$ & 0.760 & 0.795 & 0.760 & 0.779 & 0.796 & 0.361 & 0.760 \\
County FE & X & X & X & X & X & X & X \\
Year FE & X & X & X & X & X & X & X \\
\bottomrule
\end{tabular}
\begin{minipage}{\textwidth}
\vspace{0.5em}
\footnotesize \textit{Notes:} * $p < 0.1$, ** $p < 0.05$, *** $p < 0.01$. All specifications: county + year FE, state-clustered SEs. Col (5): population-weighted. Col (6): AAPI jail rate (placebo). Col (7): donut (excludes 109 counties adjacent to treated). RI $p$-value for Col (1) baseline specification (two-sided, 1,000 permutations): 0.113.
\end{minipage}
\end{table}
